\section{Relation to realization and automata theory}
The polynomial matrix upper construction realizes an entire legal
terminal density matrix, not merely one acceptance probability. In
particular, enlarging its decoder set to the full probability simplex
after tomography changes the problem. The extreme-point proof requires
the legal image of $\mathcal D_d$.

Brodsky--Pippenger \cite[Theorems~3.3 and 3.6; Lemma~3.5]{BP} characterize
bounded-error measure-once quantum language recognition by group
languages, discuss a free-group word problem, and transfer cut-point
recognition to probabilistic automata. Those statements compare languages
and cut-points; our exact profile constrains the full terminal output and
every cut of a horizon-specific stochastic machine. Ambainis--Freivalds
\cite{AF57} likewise provide an important automata-complexity comparison,
not a substitute for that legal-output quantifier.

The geometric ingredients are classical smooth homogeneous-space and
polytope arguments; compare \cite{Hall,Szarek57}. The least-rank theorem
makes the cap estimate uniform over all conditional-centroid directions,
including singular orbit types. The quaternion argument gives complete
finite-input freeness and stabilizer proofs. Bourgain--Gamburd
\cite[Theorem~1]{BG2012} is a separate deep input for the actual displayed
alphabet's noisy converse. Neither that spectral theorem nor classical
free subgroup constructions are counted as new realization theorems.

Positive realization and invariant-cone methods \cite{BF04,BR} explain
why a common positive lift, rather than unrelated one-cut factorizations,
is necessary. Hidden-process realization \cite{Heller58} concerns
consistent output processes; the terminal model is a controlled convex-valued interface. The
separate adaptive construction below specifies a nonnegative emission
instrument and a whole-transcript error bound, rather than inferring
a process law from its marginals. These
comparisons specify different mathematical questions, not an exhaustive
priority determination.

\subsection{Entropy transport and the sharp converse}
Polyanskiy--Wu \cite[Proposition~1]{PW61} prove Wasserstein continuity
of entropy for Euclidean densities under pointwise logarithmic-gradient
regularity. The entropy--transport principle itself is not new. Our
spherical mixture estimate is proved directly because compactly
supported kernels vanish and do not satisfy that regularity assumption.
The additional realization steps are the transport coupling paid for
by \emph{centroid mass} loss, the entropy defect on unions of orbitwise
small caps, and telescoping along all narrow cuts with no restriction
on intermediate widths. These give the lower power without a repeated
mixing-block logarithm.

The Hecke spectral bound \cite{LPS61,PLP61pre} is classical and deep.
It supplies the numerical norm of the actual six-rotation operator,
not a numerical experiment or a newly proved spectral theorem.
The contribution here is its combination with the stochastic
occupation argument, explicit coset profile and sparse rational upper
machine. The first rational-unitary example still imports a
non-effective Bourgain--Gamburd constant; the second rational Bloch
example has the displayed numerical width constants.

The logarithmic mismatch in the fixed-error law is thereby removed
when least-orbit and target-orbit dimensions agree. Distinct dimensions
can still give different lower and upper powers. We do not determine
optimal leading constants, the full error--horizon crossover, or rates
for every compact representation or alphabet. The stated finite-table
algorithms are polynomial in numerical horizon and inverse accuracy,
not in their binary lengths. General compact-group minimization has no
practical total-complexity claim here.

\subsection{A distinct causal approximation problem}
The process compiler constructs a nonnegative emission instrument, not
a collection of separately accurate matrix decoders. It uses a supplied
rational unit seed, rational rotations and a positive probability floor
for affine probes. The entropy comparison is performed on an auxiliary
joint trajectory with the target public marginal, so adaptive actions
are covered without conditional-error claims after rare histories.
The process algorithm is not the proof of the sharp terminal law.
That law uses the separate spectral-entropy argument and its own
assumptions; the effective spherical instance uses the LPS bound. The elementary process converse uses instead
finite statistical tests and the rational orbit-separation lattice.

\subsection{Return-free transport and varying accuracy}
Theorem~\ref{thm:driftoccupation62} uses the full action gap but no
neutral letter: the actual physical action of every filler is kept in
the transport comparison. This differs from deterministic neutral-letter
advice collapse \cite[Theorem~10]{FP}; it is not a general
stationarization theorem without synchronization. The entropy continuity
principle \cite[Proposition~1]{PW61} is classical. The additional
realization statement controls all advertised narrow cuts and prices
transport by lost centroid mass. Theorem~\ref{thm:jointlaw62} preserves
that mass-loss parameter uniformly as error vanishes; it determines the
subexponential-accuracy regime, not leading constants or every
exponential crossover. The six-command example uses the exact LPS
spectral set \cite[Theorem~1.1]{PLP61pre}, not a finite-harmonic estimate.
